\chapter{Управление проектом}

\section{Методология управления}

В ходе выполнения данного проекта была применена \textbf{гибридная методология управления}, сочетающая элементы классической водопадной модели (Waterfall) и итеративного подхода.

Выбор гибридной методологии обусловлен следующими факторами:

\begin{enumerate}
    \item \textbf{Чёткая последовательность этапов} — задачи проекта имеют выраженную зависимость: анализ предшествует проектированию, проектирование — реализации, реализация — тестированию. Это соответствует водопадной модели.
    \item \textbf{Независимость модулей} — в рамках этапов реализации (3--5) модули могут разрабатываться параллельно, что соответствует итеративному подходу.
    \item \textbf{Учебный характер проекта} — фиксированные сроки и требования не предполагают частых изменений, характерных для гибких методологий (Scrum).
\end{enumerate}

\subsection{Водопадная составляющая}

Этапы 1 и 2 (анализ и архитектура) выполнялись строго последовательно по водопадной модели:

\begin{itemize}
    \item Результаты этапа 1 (документ с требованиями) являлись неизменной базой для этапа 2.
    \item Результаты этапа 2 (схема архитектуры, структура проекта) определяли организацию кода на этапах 3--5.
    \item Каждый этап завершался формальным результатом, который утверждался перед переходом к следующему этапу.
\end{itemize}

\subsection{Итеративная составляющая}

Этапы 3--5 (реализация, алгоритмы, тестирование) выполнялись итеративно:

\begin{itemize}
    \item Каждый модуль (сбор данных, обработка, детектирование, визуализация) разрабатывался независимо.
    \item По завершении каждого модуля проводилось тестирование и интеграция с остальными компонентами.
    \item При обнаружении ошибок выполнялась доработка в рамках коротких итераций (1--2 дня).
\end{itemize}

\section{План-график работ}

Выполнение проекта было организовано в соответствии со следующим планом-графиком:

\begin{table}[H]
\centering
\caption{План-график выполнения работ}
\label{tab:schedule}
\begin{tabular}{p{1.5cm}p{4cm}p{2.5cm}p{8cm}}
\toprule
\textbf{№} & \textbf{Название этапа} & \textbf{Длит., дни} & \textbf{Результат} \\
\midrule
1 & Анализ подходов и формирование требований & 3 & Документ с анализом, функциональные и нефункциональные требования \\
\addlinespace
2 & Проектирование архитектуры & 5 & Схема архитектуры, структура проекта, диаграмма классов \\
\addlinespace
3 & Реализация сбора и потоковой обработки & 7 & Модули PacketCapture, PacketGenerator, FeatureExtractor \\
\addlinespace
4 & Разработка алгоритмов обнаружения & 7 & 3 ML-детектора, rule-based, сравнение методов \\
\addlinespace
5 & Тестирование и анализ результатов & 5 & 19 автоматизированных тестов, 8 графиков, отчёт \\
\bottomrule
\end{tabular}
\end{table}

Общая продолжительность проекта: 27 дней (включая оформление отчётной документации).

\section{Инструменты управления}

Для организации работы использовались следующие инструменты:

\begin{table}[H]
\centering
\caption{Инструменты управления проектом}
\label{tab:tools}
\begin{tabular}{p{3.5cm}p{4cm}p{8cm}}
\toprule
\textbf{Инструмент} & \textbf{Назначение} & \textbf{Описание} \\
\midrule
Git (локальный) & Контроль версий & Фиксация изменений кода, возможность отката при ошибках \\
VS Code & Среда разработки & Редактирование кода, встроенный терминал, отладка \\
pytest & Тестирование & Автоматический запуск 19 тестов, отчёт о покрытии \\
Python 3.12 & Среда исполнения & Интерпретатор языка Python \\
Markdown / \LaTeX & Документирование & Оформление отчёта \\
\bottomrule
\end{tabular}
\end{table}

\section{Оценка качества управления}

Применение гибридной методологии позволило:

\begin{enumerate}
    \item \textbf{Соблюсти сроки} — все этапы выполнены в соответствии с планом-графиком.
    \item \textbf{Обеспечить качество} — модульное тестирование каждого компонента, регулярная интеграция.
    \item \textbf{Минимизировать риски} — последовательное выполнение этапов исключило ситуацию, когда реализация начиналась без утверждённых требований.
    \item \textbf{Обеспечить гибкость} — возможность вносить изменения в отдельные модули без переработки всей системы.
\end{enumerate}

В целом, выбранная методология управления проектом показала свою эффективность для задач данного типа и масштаба.
